%! Lengths and Angles from Dot Products — Unit 1, Lesson 1.2 — Lesson Plan
\documentclass[10pt]{article}
\usepackage{linalg-article}
\usepackage{linalg-boxes}
\usepackage{graphicx}   % -article does NOT load graphicx; needed for thumbnails/screenshots
\graphicspath{{images/}}
\newcommand{\TallMath}[1]{$\displaystyle #1\rule[-1.4em]{0pt}{3.2em}$}
\newcommand{\vv}{\mathbf{v}}
\newcommand{\ww}{\mathbf{w}}
\newcommand{\UnitNumberName}{Unit 1: Vectors and Matrices \quad}
\newcommand{\LessonNumberName}{Lesson 1.2: Lengths and Angles from Dot Products}

\begin{document}
\begin{center}
    {\Large\bfseries \CourseName: \SchoolYear \\
    {\small\bfseries \UnitNumberName \LessonNumberName}}
\end{center}

\begin{tcolorbox}[colback=blush, colframe=burgundy, boxrule=0.9pt, arc=2mm,
    left=2mm, right=2mm, top=1.5mm, bottom=1.5mm]
    \textbf{Primary Objective:} Students can compute the dot product $\vv\cdot\ww$ and read it
    geometrically --- a vector's \emph{length} is $\sqrt{\vv\cdot\vv}$, the \emph{angle} between
    two vectors satisfies $\cos\theta=\dfrac{\vv\cdot\ww}{\|\vv\|\,\|\ww\|}$, and a dot product of
    \emph{zero} means the vectors are perpendicular --- justifying the sign of $\vv\cdot\ww$ as a
    measure of how much two vectors point the same way.
\end{tcolorbox}

\renewcommand{\arraystretch}{1.4}
\begin{skillbox}[Priority Ideas \& Skills]{goldbox}
\begin{tabularx}{\linewidth}{@{}X|X@{}}
\textbf{Priority Ideas \& Skills} & \textbf{Key Understandings (the ``why'')} \\[2pt]
\hline
\vspace{-6pt}
\begin{itemize}[leftmargin=1.1em, nosep, after=\vspace{2pt}]
  \item Compute a dot product $\vv\cdot\ww=v_1w_1+v_2w_2$.
  \item Recover length as $\|\vv\|=\sqrt{\vv\cdot\vv}$.
  \item Use $\cos\theta=\dfrac{\vv\cdot\ww}{\|\vv\|\,\|\ww\|}$ to find an angle.
  \item Test perpendicularity with $\vv\cdot\ww=0$.
\end{itemize}
&
\vspace{-6pt}
\begin{itemize}[leftmargin=1.1em, nosep, after=\vspace{2pt}]
  \item The dot product turns two vectors into \emph{one number} carrying length and angle.
  \item Length is just a vector's dot product with \emph{itself}, square-rooted.
  \item The \emph{sign} of $\vv\cdot\ww$ says whether vectors point the same way, cross at a right angle, or oppose.
  \item Perpendicular $\Leftrightarrow$ dot product zero --- the anchor for all of orthogonality.
\end{itemize}
\end{tabularx}
\end{skillbox}

\begin{skillbox}[{Vocabulary, Concepts, \& Theorems}]{greenbox}
\renewcommand{\arraystretch}{1.3}
\begin{tabularx}{\linewidth}{@{}l X@{}}
\textbf{Dot product} & \TallMath{\vv\cdot\ww = v_1w_1 + v_2w_2}; multiply matching components, then add. A single number (a scalar). \\
\textbf{Length (norm)} & \TallMath{\|\vv\| = \sqrt{\vv\cdot\vv} = \sqrt{v_1^2+v_2^2}}; the dot product of $\vv$ with itself, square-rooted. \\
\textbf{Unit vector} & A vector of length $1$; scale $\vv$ by $1/\|\vv\|$ to get one. \\
\textbf{Angle formula} & \TallMath{\cos\theta = \dfrac{\vv\cdot\ww}{\|\vv\|\,\|\ww\|}}; the dot product divided by the two lengths. \\
\textbf{Perpendicular (orthogonal)} & Two vectors meet at a right angle exactly when $\vv\cdot\ww=0$. \\
\end{tabularx}
\end{skillbox}

\begin{fixedskillbox}[Activate Prior Knowledge \& Spiral Review]{blush}
    The Warm-Up rehearses the Lesson~1.0 skills this lesson rebuilds into the dot product:
    \textbf{(1)} multiplying \emph{matching components} of two vectors and adding the results
    (the dot-product engine, still unnamed), \textbf{(2)} finding a length with
    $\|\vv\|=\sqrt{v_1^2+v_2^2}$, and \textbf{(3)} squaring each component and adding
    (which is exactly $\vv\cdot\vv$). Debrief item~2 out loud so the $3\text{-}4\text{-}5$
    length is fresh --- it reappears the moment we write $\|\vv\|=\sqrt{\vv\cdot\vv}$.
\end{fixedskillbox}

\begin{skillbox}[Hook]{blush}
    Project two pairs of vectors on a grid: one pair pointing \emph{almost the same way}
    ($\vv=\begin{bsmallmatrix} 3\\4 \end{bsmallmatrix}$, $\ww=\begin{bsmallmatrix} 4\\3 \end{bsmallmatrix}$)
    and one pair at a \emph{right angle}
    ($\begin{bsmallmatrix} 3\\1 \end{bsmallmatrix}$, $\begin{bsmallmatrix} -1\\3 \end{bsmallmatrix}$).
    Ask: \textbf{``Without a protractor, can one \emph{number} tell us which pair points more the
    same way?''} Let students guess, then reveal that multiplying matching components and adding
    --- the \textbf{dot product} --- gives $24$ for the first pair and $0$ for the second. Big
    number $=$ aligned; zero $=$ perpendicular. This frames the lesson: a single number that
    packages both \emph{length} and \emph{angle}.
\end{skillbox}

\begin{skillbox}[Lesson]{blush}
\begin{multicols}{2}
\textbf{1. What the dot product is.} Define $\vv\cdot\ww=v_1w_1+v_2w_2$: multiply matching
components, then add. Stress the result is a single \textbf{number}, not a vector. Compute
$\begin{bsmallmatrix} 3\\4 \end{bsmallmatrix}\cdot\begin{bsmallmatrix} 4\\3 \end{bsmallmatrix}=12+12=24$.

\textbf{2. Length is a dot product with itself.} Show $\vv\cdot\vv=v_1^2+v_2^2$, so
$\|\vv\|=\sqrt{\vv\cdot\vv}$. With $\vv=\begin{bsmallmatrix} 3\\4 \end{bsmallmatrix}$,
$\vv\cdot\vv=25$ and $\|\vv\|=5$ --- the Lesson~1.0 length formula, now seen as a dot product.
A \textbf{unit vector} scales $\vv$ by $1/\|\vv\|$.

\columnbreak

\textbf{3. The angle.} State $\vv\cdot\ww=\|\vv\|\,\|\ww\|\cos\theta$, so
$\cos\theta=\dfrac{\vv\cdot\ww}{\|\vv\|\,\|\ww\|}$. Use the pre-drawn figure and work
$\begin{bsmallmatrix} 2\\0 \end{bsmallmatrix},\begin{bsmallmatrix} 2\\2 \end{bsmallmatrix}$:
$\cos\theta=\tfrac{4}{2\cdot 2\sqrt2}=\tfrac{1}{\sqrt2}$, so $\theta=45^\circ$. Read the
\textbf{sign}: positive $\Rightarrow$ acute, zero $\Rightarrow$ right angle, negative
$\Rightarrow$ obtuse.

\textbf{4. Perpendicular $\Leftrightarrow$ dot $=0$.} Because $\cos 90^\circ=0$, two vectors are
\textbf{perpendicular} exactly when $\vv\cdot\ww=0$. Test
$\begin{bsmallmatrix} 3\\1 \end{bsmallmatrix}\cdot\begin{bsmallmatrix} -1\\3 \end{bsmallmatrix}=-3+3=0$
--- the Hook's right-angle pair.
\end{multicols}
\end{skillbox}

\begin{skillbox}[Active Monitoring]{redbox}
Circulate during the notes practice and Tier work. Watch for and correct:
\begin{itemize}[leftmargin=1.2em, nosep]
  \item writing the dot product as a \emph{vector} instead of a single number;
  \item multiplying $v_1w_2$ (crossing components) instead of matching $v_1w_1+v_2w_2$;
  \item forgetting the square root in $\|\vv\|=\sqrt{\vv\cdot\vv}$, or dividing by only one length in the cosine formula;
  \item saying ``dot product zero'' means the vectors are equal or zero, rather than \emph{perpendicular}.
\end{itemize}
Cold-call prompts: ``Is the answer a number or a vector?'' ``What is $\vv\cdot\vv$ for
$\begin{bsmallmatrix} 3\\4 \end{bsmallmatrix}$?'' ``What must the dot product be for a right angle?''
\end{skillbox}

\begin{skillbox}[Group Work \& Differentiation]{redbox}
\textbf{Are These Records Alike?} activity (tiered; mirrors the notes). Vectors are
\emph{data records} $\begin{bsmallmatrix} \text{steps} \\ \text{active min} \end{bsmallmatrix}$,
and the dot product / angle measures how similar two days' activity profiles are.
\begin{multicols}{3}
\textbf{Tier R --- Remediate.} Compute dot products and lengths of given records. Rebuilds
``multiply matching, add'' and $\|\vv\|=\sqrt{\vv\cdot\vv}$.

\columnbreak
\textbf{Tier A --- Approaching.} Find the angle between two records with the cosine formula, and
decide from the sign whether two records are alike, unrelated (perpendicular), or opposed.

\columnbreak
\textbf{Tier E --- Extension.} Build a record perpendicular to a given one, and justify why
$\vv\cdot\ww=0$ forces a right angle using the angle formula.
\end{multicols}
\end{skillbox}

\begin{skillbox}[Individual Work \& Assessment]{redbox}
\textbf{Exit Ticket} (independent, no notes): compute a dot product; find a length with
$\|\vv\|=\sqrt{\vv\cdot\vv}$; and a \textbf{perpendicularity/justification check} --- decide
whether two vectors meet at a right angle using the dot product and say what that means about the
angle. Collect and sort into ``got it / arithmetic slip / concept gap (number vs vector, or the
meaning of dot $=0$)'' to decide whether 1.3 opens with a quick re-teach.
\end{skillbox}

\begin{skillbox}[Reinforcement \& Extension]{goldbox}
\textbf{Homework:} mixed practice computing dot products and lengths, finding angles with the
cosine formula, and testing perpendicularity. \textbf{Extension:} a \emph{data-similarity}
blend --- comparing two streaming-profile records by the angle between them, previewing Unit~1's
data theme. \textbf{Preview:} Lesson~1.3, \emph{Matrices and Column Spaces} --- collecting
vectors as the columns of a matrix and asking which right-hand sides $\mathbf{b}$ the columns can
build.
\end{skillbox}
\end{document}
